% Page 25 — Alban – XIVe siècle

\begin{center}
  {\Large\itshape Alban – XIV\textsuperscript{ème} siècle}
\end{center}
\vspace{1.2em}

\noindent
\lettrine{R}{écemment} à 2,5 kilomètres d’Alban a été apposée une plaque invitant d’éventuels touristes à faire un tour au « Roc des Anglais ». Qu’ont à faire avec cette masse de rochers qui surplombe les quelques maison de Nougaret nos amis britanniques ? La réponse fait voisiner l’histoire et la légende. C’est dans ces parages qu’a bien pu se situer un village de même nom dont le château dit de « Saint Michel de Nougaret » est cité dans plusieurs actes du XIII\textsuperscript{ème} siècle débutant. Si le touriste est un sportif, il pourra, par un sentier de chèvres, atteindre une sorte de réservoir rocheux, « la paissiayro » – citerne (peut-être) pour alimenter en eau les maisonnettes abritées dans les flancs du rocher taillé à angles droits. Le professeur Soutou de Toulouse en a situé une vingtaine. Le visiteur agile ne craignant pas la caresse des ronces y pourra faire des prouesses de cabri pour passer d’un point de repère à un autre et découvrir par lui-même les vestiges d’un lointain passé. La tradition veut que ce coin de montagne ait été le témoin de bagarres entre Français et les sujets de « Sa Gracieuse Majesté ». Les Anglais n’avaient pas attendu le traité de Brétigny (1360) pour occuper un tiers de la France. C’est la fameuse guerre de Cent ans. Le 25 janvier 1309, ils se sont emparés du château de Roquecézière. Leurs incursions ont pu les amener tout près de chez nous. Il ne semble pas que ces insulaires aient assiégé Alban. Mais ils ont bien pu attaquer le village de Nougaret et l’occuper. Il y aurait eu du côté de Tonguière un pré qui, depuis, se serait appelé le « pré des mille hommes » à cause d’une concentration de troupes ennemies en ce lieu.
